% Source: https://artofproblemsolving.com/community/c3h1413242_2m52n1_can_be_written_as_a_sum_of_two_perfect_squar
% Stable identity: aops-topic-1413242-q01
Since $n\ge1$, the exponent $2^n+1$ is odd; write $2^n+1=2j+1$. Then
\[
5^{2^n+1}=5\cdot25^{j}=(1+4)\cdot\left(5^{j}\right)^2=\left(5^{j}\right)^2+\left(2\cdot5^{j}\right)^2 .
\]
\textbf{If $m=2t$ is even:}
\[
2^m\cdot5^{2^n+1}=\left(2^t5^j\right)^2+\left(2^{t+1}5^j\right)^2 .
\]
\textbf{If $m=2t+1$ is odd:} use $2(x^2+y^2)=(x+y)^2+(x-y)^2$ with $x=5^j$, $y=2\cdot5^j$:
\[
2^m\cdot5^{2^n+1}=4^t\cdot2\left(x^2+y^2\right)=4^t\left[(3\cdot5^j)^2+(5^j)^2\right]=\left(3\cdot2^t5^j\right)^2+\left(2^t5^j\right)^2 .
\]
In both cases the number is a sum of two squares of positive integers. $\blacksquare$

(For example $m=1$, $n=1$: $2\cdot125=250=15^2+5^2$.)
